\section{System}
\label{sec:system}

\subsection{Roles}

\begin{table}[H]
\centering
\begin{tabularx}{\textwidth}{L{0.26\textwidth} Y}
\toprule
\textbf{Role} & \textbf{Daemons} \\
\midrule
Entry node & \kn{master}, \kn{collector}, \kn{negotiator}, \kn{schedd}, \kn{startd} (limited), \kn{shared\_port} \\
Worker & \kn{master}, \kn{startd} \\
User desktop & Worker daemons plus the submit client; \textbf{no schedd} \\
Submit-only client & No daemons at all. Typically laptops \\
\bottomrule
\end{tabularx}
\caption{Daemon layout by role.}
\end{table}

There is one schedd, on the entry node. A schedd on a laptop dies with the lid
and takes its queue with it, so centralising it means a user's jobs survive
their own machine sleeping, which is the normal state of a laptop between two
meetings.

Donation is expected of desktops, which sit still and stay plugged in. A laptop
that only submits is a normal member of the pool rather than an opt-out, and the
reason is a protocol asymmetry rather than a preference: a worker must be
reachable \emph{inbound} to be claimed, so an off-network laptop registers,
matches, and then fails the claim, while a submit client makes only
\emph{outbound} connections and therefore works through a tunnel and from
outside the office.

\subsection{Choosing the entry node}

The entry node is a low-power four-core machine rather than the pool's
highest-core one. The role is bounded by \kn{fsync} latency on the job queue,
RAM per shadow and NIC bandwidth, not by CPU, so sixteen threads buy nothing
there. The four-core class is also the fleet's repeated unit, which makes it the
honest thing to test on: a production entry node will be commodity hardware, and
testing on the one scarce machine would teach nothing transferable. The scaling
path is horizontal, by adding schedds or splitting the collector and negotiator
off, never vertical.

\subsection{Donate everything, protect dynamically}

The instinct when pooling somebody's desktop is to hold capacity back. This
design does the opposite on workers, because static reservation and the
contention policy of \S\ref{sec:admission} defend against the same thing.
Holding capacity back only idles it in the common case, which is that nobody is
using the machine, while the policy would suspend anyway the moment someone
arrives.

The entry node is deliberately different. It carries the collector, negotiator
and schedd, and if the schedd cannot fork shadows or \kn{fsync} its queue then
nothing anywhere in the pool runs, so it reserves about half its cores and
roughly two thirds of its memory.

\subsection{Advertisement cadence}

\kn{UPDATE\_INTERVAL} defaults to 300 seconds, a value sized for pools of
thousands. At this scale an advertisement every 30 seconds is free, and the
default costs twice: \kn{condor\_status} showed a frozen value for two and a half
minutes while load was driven from 1 to 15, and after a central-manager restart
workers took roughly twice the interval to reappear, one cycle to discover their
cached security session was dead and another to re-authenticate.

Lowering it does not change how fast policy reacts. The startd evaluates its
expressions every \kn{POLLING\_INTERVAL}, five seconds, regardless. The default
made the policy invisible and re-registration slow. It did not make the policy
slow.

\subsection{Workload constraints}

Eviction is the normal case in this design rather than an exception, and users
have to be told so. Jobs should be many and short rather than few and long; they
must be idempotent, restartable from scratch without corrupting state; long work
should self-checkpoint. ``My job keeps restarting'' is a property of a desktop
pool, not a bug.
